\section{Method}
\label{sec:method}

\subsection{Simulation reservoir}

The reservoir has six qubits and two shallow data reuploading layers
~\cite{perezsalinas2020data}. Each layer encodes the scalar input on every
qubit with heterogeneous $R_y$ and $R_z$ rotations, applies neighboring
$R_{zz}$ couplings, and mixes locally with frozen non-Clifford angles. The
circuit parameters are fixed before forecasting, and additional frozen
realizations quantify sensitivity. One computational basis measurement yields
six local $Z_i$ signs and five neighboring products $Z_iZ_{i+1}$ at no
additional quantum cost.

Ideal statevector simulation is performed before forecasting. All 11 features
change with the input and exceed the specified nonlinearity threshold. They
all contribute distinct numerical information. Because several respond in
similar ways, their effective number is 5.15. A leak of 0.2
and ten causal lags provide temporal context. Ridge regression produces the
next step forecast, and performance is measured by test NRMSE. These checks
show that the reservoir supplies a varied, nonlinear feature set. They do not evaluate
quantum advantage.

\subsection{Adaptive shot aggregation}

Let $x_t^{(n)}$ denote the processed features from shot index $n$ at time
$t$. For a divisor $k$ of $N$, the $N$ shots are partitioned into
$G=N/k$ disjoint groups, producing
\begin{equation}
  x_t^{(g)}=\frac{1}{k}\sum_{n\in g}x_t^{(n)},
  \qquad g=1,\ldots,G .
  \label{eq:grouping}
\end{equation}
Each grouped sample shares the same target. At evaluation, the $G$
predictions for one step are averaged. This equals prediction from their mean
features, so any difference between EV and Grouped is created during
training rather than evaluation.

Measurement indices do not identify physical trajectories. A random
permutation is therefore drawn at every time step before groups are formed, and
the analysis is repeated across random grouping assignments. A fixed assignment across time
is included as a sensitivity analysis.

\subsection{Fixed execution budget and training}

The quantum cost is $TN$ circuit executions. Three specified settings use
exactly $10{,}000$ executions. They allocate $(T,N)=(250,40)$, $(500,20)$, and
$(1000,10)$ to the number of time steps and shots. All divisors of $N$ are
considered, including Raw at $k=1$ and EV at $k=N$. Three expanding
train-and-check splits preserve time order: each candidate pair $(k,\lambda)$
is fitted on an earlier block and scored on the following block. The best pair
is then fixed, the readout is refitted to all data preceding the test block,
and the final 20\% is evaluated once. Normalization and feature scaling use
training data only.
